\documentclass[../DoAn.tex]{subfiles}

\begin{document}

\section{Thiết kế thực nghiệm tổng thể}
\label{section:experiment-design}

Chương này đánh giá hệ thống theo hai luồng đã trình bày ở Chương 4 và
Chương 5: trích xuất dữ liệu từ PDF và khai thác dữ liệu đã trích xuất để truy
xuất, hỏi đáp. Các thí nghiệm được tổ chức theo câu hỏi nghiên cứu thay vì lặp
lại mô tả kiến trúc.

\begin{table}[H]
\centering
\small
\renewcommand{\arraystretch}{1.2}
\caption{Câu hỏi nghiên cứu và kết quả dùng để trả lời}
\label{tab:experiment-rq}
\begin{tabularx}{\textwidth}{
    |>{\raggedright\arraybackslash}p{0.10\textwidth}
    |>{\raggedright\arraybackslash}X
    |>{\raggedright\arraybackslash}p{0.24\textwidth}
    |>{\raggedright\arraybackslash}p{0.22\textwidth}|
}
\hline
\textbf{Mã} & \textbf{Câu hỏi nghiên cứu} & \textbf{Tầng đánh giá} &
\textbf{Độ đo chính} \\
\hline
RQ1 &
Luồng ingest chính với probe, định tuyến vùng, OCR, text extractor và table
extractor tạo được dữ liệu hợp lệ đến mức nào? &
Ingest chính và định tuyến vùng &
Tỷ lệ thành công, route trace, token F1, metadata, thời gian xử lý \\
\hline
RQ2 &
Hybrid TATR có cải thiện tái dựng cấu trúc và nội dung bảng không? &
Trích xuất bảng &
Cell F1, Structure F1, Exact CSV \\
\hline
RQ3 &
Truy xuất kết hợp có cải thiện khả năng tìm đúng bằng chứng so với BM25 và
truy xuất ngữ nghĩa riêng lẻ không? &
Truy xuất thông tin &
Hit@k, Recall@k, MRR, NDCG \\
\hline
RQ4 &
Kiểm tra bằng chứng và dẫn chứng nguồn có giúp kiểm soát câu trả lời thiếu căn
cứ không? &
Hỏi đáp đầu cuối &
Answer Match, Evidence Match, Grounded Rate, Hallucination Rate \\
\hline
RQ5 &
Chia đoạn có nhận thức bảng có giúp truy xuất đúng hàng và ô không? &
Truy xuất bảng &
Table, Row, Column và Cell Match@k \\
\hline
\end{tabularx}
\end{table}

\begin{figure}[H]
\centering
\small
\begin{tabular}{c}
\fbox{\parbox{0.82\linewidth}{\centering
RQ1--RQ2: PDF \(\rightarrow\) vùng nội dung \(\rightarrow\) block bảng/văn bản}}\\[0.15cm]
\(\downarrow\)\\[-0.05cm]
\fbox{\parbox{0.82\linewidth}{\centering
RQ3--RQ5: chunk và chỉ mục \(\rightarrow\) bằng chứng \(\rightarrow\) câu trả lời có dẫn chứng}}\\[0.15cm]
\(\downarrow\)\\[-0.05cm]
\fbox{\parbox{0.82\linewidth}{\centering
Phân tích loại bỏ thành phần \(\rightarrow\) phân tích lỗi \(\rightarrow\) giới hạn kết quả}}
\end{tabular}
\caption{Bản đồ các nhóm thực nghiệm trong chương}
\label{fig:experiment-map}
\end{figure}

\section{Thiết lập chung}
\label{section:experiment-setup}

\subsection{Phạm vi các tập đánh giá}

Các tập đánh giá được chia thành ba nhóm. QCDT được dùng để đánh giá luồng
hỏi đáp chính và tác động đầu cuối của ingest. SciFact và QASPER được dùng như
tập phụ để khảo sát khả năng chuyển miền sang tài liệu khoa học. PubTables-1M
và các tập OCR được dùng để đánh giá các backend bảng/OCR được pipeline ingest
gọi trực tiếp \cite{smock2022pubtables}. Các kết quả phát hiện bố cục bằng
backend \texttt{model\_layout} tùy chọn không được dùng làm kết luận chính vì
pipeline chính của đồ án là probe--region routing--OCR/text/table.

\begin{table}[H]
\centering
\small
\renewcommand{\arraystretch}{1.2}
\caption{Vai trò của các tập dữ liệu trong thực nghiệm}
\label{tab:experiment-datasets}
\begin{tabularx}{\textwidth}{
    |>{\raggedright\arraybackslash}p{0.21\textwidth}
    |>{\raggedright\arraybackslash}p{0.19\textwidth}
    |>{\raggedright\arraybackslash}X
    |>{\raggedright\arraybackslash}p{0.20\textwidth}|
}
\hline
\textbf{Tập dữ liệu} & \textbf{Vai trò} & \textbf{Nội dung đánh giá} &
\textbf{Phạm vi diễn giải} \\
\hline
QCDT &
Chính &
40 câu hỏi trên tài liệu Quy chế đào tạo của Đại học Bách khoa Hà Nội
\cite{hust2025trainingregulation}. &
Tài liệu quy định tiếng Việt có lớp văn bản. \\
\hline
SciFact &
Phụ &
300 câu hỏi dùng để đánh giá truy xuất và hỏi đáp khoa học
\cite{wadden2020scifact}. &
Phân tích khả năng chuyển miền. \\
\hline
QASPER &
Phụ &
100 câu hỏi dùng để đánh giá truy xuất evidence trên các đoạn trong bài báo;
500 câu hỏi dùng để đánh giá hỏi đáp đầu cuối trên bài báo khoa học
\cite{dasigi2021qasper}. &
Tập dữ liệu gồm câu hỏi tự nhiên, đáp án và bằng chứng gắn với bài báo khoa học;
dùng để phân tích giới hạn với câu hỏi dài và nhiều bằng chứng. \\
\hline
PubTables-1M &
Thành phần &
Đánh giá tái dựng cấu trúc bảng trên tập PubTables structure, gồm hàng, cột, ô,
gán văn bản vào ô và so khớp CSV. &
Tập 25 mẫu dùng để so sánh backend bảng; tập 100 mẫu proxy dùng để kiểm tra
Hybrid TATR trên số lượng mẫu lớn hơn và phân tích lỗi theo nhóm bảng. \\
\hline
OCR scan, FUNSD &
Thành phần &
Đánh giá backend OCR trên tập scan tổng hợp và tập biểu mẫu FUNSD đã chuyển
về định dạng benchmark OCR. &
Dùng để kiểm tra khả năng nhận dạng chữ của nhánh OCR; không thay thế đánh
giá đầy đủ PDF scan tiếng Việt thực tế. \\
\hline
\end{tabularx}
\end{table}

\subsection{Quy ước độ đo}

Các công thức Success Rate, Token F1, CER, Latency, Hit@k, Recall@k, MRR,
NDCG, IoU, Precision, Recall, F1, Exact CSV và Match@k ở mức bảng đã được
định nghĩa tại Chương 3. Trong chương này, \(k\) được ghi cùng từng tập dữ
liệu. Ký hiệu \(k\) là số kết quả đứng đầu trong danh sách truy xuất được dùng
để tính độ đo. Ví dụ, Hit@10 kiểm tra trong 10 kết quả đầu tiên có bằng chứng
đúng hay không. SciFact dùng \(k=5\), QCDT dùng \(k=10\), QASPER dùng
\(k=20\), còn thực nghiệm truy xuất bảng báo cáo cả \(k=5\) và \(k=10\).

Answer Match đo tỷ lệ câu trả lời khớp theo quy tắc của bộ đánh giá. Evidence
Match đo tỷ lệ tìm hoặc trích dẫn được ít nhất một bằng chứng đúng. Grounded
Rate đo tỷ lệ câu trả lời có dẫn chứng hợp lệ theo quy ước triển khai.
Hallucination Rate đo tỷ lệ câu trả lời bị bộ đánh giá xác định là thiếu căn cứ.
Các độ đo này không chứng minh câu trả lời đúng về mặt ngữ nghĩa trong mọi
trường hợp; vì vậy chúng được đọc cùng Answer Match và phần phân tích lỗi ở
Mục~\ref{section:error-analysis}.

\subsection{Cấu hình so sánh}

\begin{table}[H]
\centering
\small
\renewcommand{\arraystretch}{1.2}
\caption{Các cấu hình chính trong thực nghiệm}
\label{tab:experiment-configurations}
\begin{tabularx}{\textwidth}{
    |>{\raggedright\arraybackslash}p{0.13\textwidth}
    |>{\raggedright\arraybackslash}p{0.28\textwidth}
    |>{\raggedright\arraybackslash}X|
}
\hline
\textbf{Mã} & \textbf{Cấu hình} & \textbf{Mục đích} \\
\hline
C0 &
\texttt{routed\_grounded} &
Cấu hình hỏi đáp chính gồm định tuyến câu hỏi, truy xuất, kiểm tra bằng chứng,
sinh câu trả lời và dẫn chứng. \\
\hline
C1 &
BM25 / dense / hybrid &
So sánh các tín hiệu truy xuất \cite{robertson2009bm25,
reimers2019sentencebert,cormack2009rrf}. \\
\hline
C2 &
Default / TATR / Hybrid TATR &
So sánh các backend tái dựng bảng. \\
\hline
C3 &
Normal / table-aware / table-aware cell &
Đo tác động của đơn vị truy xuất bảng, hàng và ô. \\
\hline
C4 &
\texttt{no\_router}, \texttt{no\_evidence\_checker},
\texttt{no\_citation\_grounding} &
Đo vai trò của từng thành phần trong luồng hỏi đáp. \\
\hline
C5 &
Region OFF / Region ON &
Đo tác động của \texttt{region\_routed} trong pipeline ingest chính bằng cách
chạy cùng một tập PDF với và không có định tuyến vùng. \\
\hline
\end{tabularx}
\end{table}

Kết quả được lấy trực tiếp từ các tệp JSON trong thư mục \texttt{results}. Chương
không sử dụng cấu hình tự động thử lại tuyến truy xuất. Hai trạng thái
\texttt{expand\_retrieval} và \texttt{switch\_strategy}, nếu xuất hiện, chỉ là
trạng thái đánh giá bằng chứng trong luồng chuẩn.

\section{Thực nghiệm theo tầng pipeline}
\label{section:layer-experiments}

\subsection{Probe, định tuyến và trích xuất trong luồng ingest chính}
\label{subsection:ingest-results}

Pipeline ingest chính được đánh giá theo đúng luồng
probe--region routing--OCR/text/table. Mục tiêu của phần này không phải chứng
minh một mô hình layout riêng, mà kiểm tra bốn điểm: probe chọn được hướng xử
lý ban đầu, region routing tạo được vết định tuyến, các backend OCR/text/table tạo
được block hợp lệ, và cơ chế fallback không làm pipeline thất bại khi backend đầu
tiên không phù hợp. Kết quả được tổng hợp từ các tệp trong thư mục
\texttt{results} bằng script \texttt{scripts/benchmark\_main\_ingest\_pipeline.py}.

\begin{table}[H]
\centering
\scriptsize
\renewcommand{\arraystretch}{1.2}
\caption{Đánh giá luồng ingest chính theo probe, định tuyến và backend}
\label{tab:ingest-component-results}
\begin{tabularx}{\textwidth}{
    |>{\raggedright\arraybackslash}p{0.21\textwidth}
    |>{\raggedright\arraybackslash}p{0.18\textwidth}
    |>{\raggedright\arraybackslash}p{0.20\textwidth}
    |>{\raggedright\arraybackslash}X|
}
\hline
\textbf{Nội dung kiểm tra} & \textbf{Dữ liệu} & \textbf{Độ đo} &
\textbf{Kết quả và diễn giải} \\
\hline
Probe và chọn backend cấp tài liệu &
31 PDF gồm QCDT, PDF bảng và scan tổng hợp &
Probe accuracy, backend routing accuracy &
Probe chọn đúng hướng xử lý cho 31/31 tài liệu, đạt accuracy 1,000. Backend
cấp tài liệu sau probe/fallback cũng đúng 31/31: 30 tài liệu đi qua OCR và tài
liệu QCDT đi qua \texttt{region\_routed} khi bật định tuyến vùng. \\
\hline
Định tuyến vùng trên QCDT &
Một tài liệu QCDT có text layer và bảng &
Route expectation, route trace, route backend &
Kỳ vọng QCDT cần cả nhánh text và nhánh bảng; hệ thống đạt 1/1 ở mức tài liệu.
Bật region routing tạo 9 block bảng và 494 block có vết định tuyến. Trong đó
484 block đi qua nhánh text, 8 block qua nhánh table, 1 block qua Hybrid TATR
và 1 block placeholder. \\
\hline
Backend OCR &
25 mẫu scan tổng hợp và 25 mẫu FUNSD OCR &
Token F1, CER &
Scan tổng hợp đạt Token F1 1,000 và CER 0,000. FUNSD đạt Token F1 0,826 và
CER 0,466, cho thấy OCR giảm chất lượng rõ rệt trên dữ liệu nhiễu hơn
\cite{du2020ppocr}. \\
\hline
Backend bảng &
100 mẫu PubTables-1M proxy &
Cell F1, Structure F1, Text Assignment F1, Exact CSV &
Hybrid TATR đạt Cell F1@0,50 là 0,947, Structure F1 là 0,778,
Text Assignment F1 là 0,997 và Exact CSV là 0,530. Đây là đánh giá backend
bảng được pipeline gọi khi gặp vùng bảng, không phải đánh giá OCR đầu cuối. \\
\hline
Đầu ra block/chunk/metadata &
31 PDF ở thí nghiệm Region OFF/ON &
Block success rate, metadata completeness &
Region ON đạt block/chunk success rate 1,000, tạo trung bình 64,45 block,
53,23 chunk và 15,94 block có route trace trên mỗi tài liệu. Trên các block mẫu,
tỷ lệ có page id, bbox, block type, text và route trace đều đạt 1,000. \\
\hline
Fallback và chi phí xử lý &
31 PDF ở thí nghiệm Region OFF/ON &
Fallback count, fallback success rate, latency &
Fallback được kích hoạt ở 30 tài liệu và thành công 30/30, không có tài liệu lỗi.
Đổi lại, latency trung bình tăng từ 6,189 s lên 16,556 s khi bật region routing. \\
\hline
\end{tabularx}
\end{table}

Bảng~\ref{tab:ingest-component-results} không báo F1 phát hiện bbox bố cục vì
thí nghiệm này không dùng nhãn vùng thủ công cho region routing. Các giá trị
accuracy trong bảng được hiểu ở mức tài liệu hoặc mức kỳ vọng route của tài
liệu QCDT, không phải Precision/Recall/F1 cho từng bbox vùng. Nếu muốn đo F1
theo từng loại vùng, cần bổ sung một tập PDF thật có nhãn chuẩn cho vùng text,
bảng, ảnh và caption. Trong phạm vi đồ án, chất lượng ingest chính được đọc qua
success rate, backend được chọn, số block bảng, vết định tuyến, metadata đầu ra,
OCR, backend bảng và ảnh hưởng đầu cuối ở Mục~\ref{subsection:region-routing-results}.
Kết quả OCR 1,000 trên scan tổng hợp chỉ được xem là kiểm tra chức năng trên dữ
liệu sạch, không dùng để khẳng định hệ thống đã xử lý hoàn chỉnh PDF scan tiếng
Việt trong mọi trường hợp.

\subsection{Định tuyến theo vùng}
\label{subsection:region-routing-results}

Thực nghiệm theo cặp gồm 31 PDF: một tài liệu QCDT, 25 PDF đại diện cho bảng
và 5 PDF scan tổng hợp. Trên 30 PDF scan hoặc đại diện, cả hai cấu hình chủ yếu
chuyển sang OCR theo cơ chế dự phòng.
Trên QCDT, bật định tuyến vùng làm backend thay đổi từ \texttt{text} sang
\texttt{region\_routed}.

\begin{table}[H]
\centering
\small
\renewcommand{\arraystretch}{1.2}
\caption{So sánh định tuyến vùng trên 31 PDF}
\label{tab:region-routing-results}
\begin{tabularx}{\textwidth}{
    |>{\raggedright\arraybackslash}p{0.26\textwidth}
    |>{\centering\arraybackslash}p{0.19\textwidth}
    |>{\centering\arraybackslash}p{0.19\textwidth}
    |>{\raggedright\arraybackslash}X|
}
\hline
\textbf{Độ đo} & \textbf{Tắt region} & \textbf{Bật region} &
\textbf{Chênh lệch} \\
\hline
Success rate & 1,000 & 1,000 & Không đổi \\
\hline
Số block bảng trung bình & 0,97 & 1,26 & \(+0,29\) \\
\hline
Block có route trace trung bình & 0,00 & 15,94 & \(+15,94\) \\
\hline
Thời gian trung bình & 6,189 s & 16,556 s & \(+10,367\) s \\
\hline
\end{tabularx}
\end{table}

Riêng QCDT, cấu hình tắt region tạo 0 block bảng, còn cấu hình bật region tạo
9 block bảng và gắn route trace cho 494 block. Đổi lại, thời gian tăng từ
0,286 s lên 70,189 s trong lần chạy cold-run có tải backend bảng. Kết quả này
chứng minh vai trò điều phối và truy vết của region routing trên QCDT, nhưng
chưa chứng minh định tuyến vùng luôn cải thiện chất lượng hoặc tốc độ trên mọi
PDF.

Để kiểm tra ảnh hưởng đầu cuối, đồ án chạy thêm một thí nghiệm trên cùng 40
câu hỏi QCDT: tắt region rồi ingest--index--retrieval--QA, sau đó bật region và
chạy lại toàn bộ luồng. Cả hai cấu hình dùng index BM25 để cô lập tác động của
ingest và region routing.

\begin{table}[H]
\centering
\small
\renewcommand{\arraystretch}{1.2}
\caption{Thực nghiệm đầu cuối region OFF/ON trên QCDT}
\label{tab:region-routing-e2e-results}
\begin{tabularx}{\textwidth}{
    |>{\raggedright\arraybackslash}p{0.21\textwidth}
    |>{\centering\arraybackslash}p{0.12\textwidth}
    |>{\centering\arraybackslash}p{0.12\textwidth}
    |>{\centering\arraybackslash}p{0.12\textwidth}
    |>{\centering\arraybackslash}p{0.12\textwidth}
    |>{\centering\arraybackslash}p{0.12\textwidth}
    |>{\centering\arraybackslash}X|
}
\hline
\textbf{Cấu hình} & \textbf{Chunk} & \textbf{Bảng} & \textbf{Hit@10} &
\textbf{Answer Match} & \textbf{Evidence Match} & \textbf{QA latency} \\
\hline
Region OFF & 246 & 0 & 0,550 & 0,775 & 1,000 & 73,032 ms \\
\hline
Region ON & 326 & 9 & 0,475 & 0,700 & 1,000 & 73,515 ms \\
\hline
\end{tabularx}
\end{table}

Kết quả đầu cuối cho thấy region routing làm giàu dữ liệu bảng và tăng số chunk,
nhưng không tự động cải thiện truy xuất hoặc hỏi đáp trên mọi câu hỏi QCDT.
Trong cấu hình BM25-only, nhiều chunk bảng mới làm không gian truy xuất rộng
hơn, khiến Hit@10 và Answer Match giảm nhẹ. Vì vậy, region routing nên được
diễn giải là tầng trích xuất và truy vết nguồn; để chuyển lợi ích bảng thành chất
lượng QA cần thêm table-aware retrieval hoặc reranking phù hợp. Thí nghiệm này
cố ý giữ index BM25-only để cô lập tác động của ingest; vì vậy nó chưa phủ hết
các cấu hình truy vấn có ưu tiên metadata bảng.

\subsection{Tái dựng bảng}
\label{subsection:table-reconstruction-results}

Ba backend được chạy trên cùng 25 mẫu PubTables-1M. Hybrid TATR trong thực
nghiệm này kết hợp cấu trúc hình học từ Table Transformer với hộp từ lấy từ lớp
văn bản PDF; OCR không được gọi trực tiếp bên trong module Hybrid TATR
\cite{smock2022pubtables}.

\begin{table}[H]
\centering
\scriptsize
\renewcommand{\arraystretch}{1.15}
\caption{Kết quả tái dựng bảng trên 25 mẫu PubTables-1M}
\label{tab:table-reconstruction-results}
\begin{tabularx}{\textwidth}{
    |>{\raggedright\arraybackslash}p{0.15\textwidth}
    |>{\centering\arraybackslash}p{0.10\textwidth}
    |>{\centering\arraybackslash}p{0.10\textwidth}
    |>{\centering\arraybackslash}p{0.10\textwidth}
    |>{\centering\arraybackslash}p{0.11\textwidth}
    |>{\centering\arraybackslash}p{0.13\textwidth}
    |>{\centering\arraybackslash}X|
}
\hline
\textbf{Backend} & \textbf{Detect F1} & \textbf{Cell F1@0,50} &
\textbf{Cell F1@0,75} & \textbf{Structure F1} &
\textbf{Text Assignment F1} & \textbf{Exact CSV} \\
\hline
Default & 0,940 & 0,650 & 0,149 & 0,199 & 0,909 & 0,000 \\
\hline
TATR & 0,987 & 0,491 & 0,103 & 0,010 & 0,015 & 0,000 \\
\hline
\textbf{Hybrid TATR} & \textbf{0,987} & \textbf{0,958} & \textbf{0,944} &
\textbf{0,772} & \textbf{0,999} & \textbf{0,480} \\
\hline
\end{tabularx}
\end{table}

Hybrid TATR cải thiện rõ Cell F1, Structure F1 và khả năng gán chữ vào ô.
Tuy nhiên, Exact CSV chỉ đạt 0,480. Các lỗi còn lại tập trung ở ô gộp, tiêu đề
nhiều hàng và sai khác nhỏ trong thứ tự hàng--cột. Vì vậy, kết quả trả lời RQ2
theo hướng tích cực nhưng không cho phép kết luận bài toán tái dựng bảng đã được
giải quyết hoàn toàn.

Để kiểm tra kết luận này trên tập lớn hơn, đồ án chạy thêm 100 mẫu
PubTables-1M proxy. Thực nghiệm mở rộng tập trung vào TATR-only và Hybrid
TATR, vì mục tiêu là kiểm tra vai trò của việc gán hộp từ vào cấu trúc hình học.

\begin{table}[H]
\centering
\scriptsize
\renewcommand{\arraystretch}{1.15}
\caption{Kết quả tái dựng bảng mở rộng trên 100 mẫu PubTables-1M}
\label{tab:table-reconstruction-100-results}
\begin{tabularx}{\textwidth}{
    |>{\raggedright\arraybackslash}p{0.18\textwidth}
    |>{\centering\arraybackslash}p{0.11\textwidth}
    |>{\centering\arraybackslash}p{0.11\textwidth}
    |>{\centering\arraybackslash}p{0.12\textwidth}
    |>{\centering\arraybackslash}p{0.14\textwidth}
    |>{\centering\arraybackslash}p{0.12\textwidth}
    |>{\centering\arraybackslash}X|
}
\hline
\textbf{Backend} & \textbf{Detect F1} & \textbf{Cell F1@0,50} &
\textbf{Structure F1} & \textbf{Text Assignment F1} & \textbf{Exact CSV} &
\textbf{Latency} \\
\hline
TATR-only & 0,980 & 0,505 & 0,021 & 0,032 & 0,000 & 0,304 s \\
\hline
\textbf{Hybrid TATR} & \textbf{0,980} & \textbf{0,947} & \textbf{0,778} &
\textbf{0,997} & \textbf{0,530} & \textbf{0,268 s} \\
\hline
\end{tabularx}
\end{table}

Trên 100 mẫu, Hybrid TATR tiếp tục cải thiện mạnh so với TATR-only ở các độ đo
liên quan đến nội dung bảng. Exact CSV tăng từ 0,480 trên tập 25 mẫu lên 0,530
trên tập 100 mẫu, nhưng vẫn thấp hơn nhiều so với Cell F1 và Text Assignment F1.
Điều này cho thấy hệ thống thường gán đúng phần lớn ô và văn bản, nhưng chỉ cần
một sai khác nhỏ trong cấu trúc hoặc thứ tự ô là toàn bộ CSV bị tính sai. Cần lưu
ý rằng tập 100 mẫu này là tập proxy dùng hộp từ chuẩn hóa của PubTables-1M,
không phải đánh giá OCR đầu cuối trên tài liệu scan thực tế.

\begin{table}[H]
\centering
\scriptsize
\renewcommand{\arraystretch}{1.15}
\caption{Phân tích Hybrid TATR theo nhóm bảng trên 100 mẫu}
\label{tab:hybrid-tatr-group-results}
\begin{tabularx}{\textwidth}{
    |>{\raggedright\arraybackslash}p{0.24\textwidth}
    |>{\centering\arraybackslash}p{0.08\textwidth}
    |>{\centering\arraybackslash}p{0.12\textwidth}
    |>{\centering\arraybackslash}p{0.12\textwidth}
    |>{\centering\arraybackslash}p{0.13\textwidth}
    |>{\centering\arraybackslash}X|
}
\hline
\textbf{Nhóm bảng} & \textbf{Số mẫu} & \textbf{Exact CSV} &
\textbf{Cell F1@0,50} & \textbf{Structure F1} & \textbf{Text Assignment F1} \\
\hline
Bảng đơn giản & 4 & 0,750 & 0,795 & 0,795 & 1,000 \\
\hline
Bảng có ô gộp & 48 & 0,229 & 0,942 & 0,630 & 0,998 \\
\hline
Tiêu đề nhiều hàng & 26 & 0,385 & 0,963 & 0,610 & 0,998 \\
\hline
Bảng không có đường kẻ & 88 & 0,545 & 0,951 & 0,773 & 0,997 \\
\hline
\end{tabularx}
\end{table}

Các nhóm trong Bảng~\ref{tab:hybrid-tatr-group-results} có thể chồng lấp, ví dụ
một bảng vừa có ô gộp vừa không có đường kẻ. Nhãn ``không có đường kẻ'' được
xác định bằng heuristic trên ảnh bảng, không phải nhãn chính thức của PubTables.
Kết quả cho thấy lỗi Exact CSV tập trung nhiều ở bảng có ô gộp và tiêu đề nhiều
hàng; đây là hai nhóm cần cải thiện nếu muốn tái dựng bảng chính xác tuyệt đối.
Do các nhóm chồng lấp, tổng số mẫu trong các nhóm không cộng đúng bằng 100.

\subsection{Chia đoạn có nhận thức bảng}
\label{subsection:table-aware-results}

Thực nghiệm mở rộng gồm 39 câu hỏi bảng trên 6 bảng của tài liệu QCDT, dùng
BM25 trên ba corpus được tạo từ cùng dữ liệu ingest. Các câu hỏi bao gồm tra ô
trực tiếp, tra ngược theo giá trị, tra theo khoảng điều kiện và một nhóm bảng có
header rộng. Kết quả top-5 là kết quả chính; top-10 được báo cáo để kiểm tra liệu
bằng chứng ô có tồn tại ở thứ hạng thấp hơn.

\begin{table}[H]
\centering
\small
\renewcommand{\arraystretch}{1.2}
\caption{Ảnh hưởng của chia đoạn có nhận thức bảng}
\label{tab:table-aware-results}
\begin{tabularx}{\textwidth}{
    |>{\raggedright\arraybackslash}p{0.31\textwidth}
    |>{\centering\arraybackslash}p{0.14\textwidth}
    |>{\centering\arraybackslash}p{0.14\textwidth}
    |>{\centering\arraybackslash}p{0.14\textwidth}
    |>{\centering\arraybackslash}X|
}
\hline
\textbf{Cấu hình} & \textbf{Table@5} & \textbf{Row@5} &
\textbf{Column@5} & \textbf{Cell@5} \\
\hline
Chunk thông thường & 1,000 & 0,000 & 1,000 & 0,000 \\
\hline
Table-aware, không có cell chunk & 1,000 & 0,974 & 1,000 & 0,000 \\
\hline
Table-aware và cell chunk & 1,000 & 1,000 & 1,000 & 0,487 \\
\hline
\end{tabularx}
\end{table}

Khi tăng lên top-10, cấu hình có cell chunk đạt Cell Match 0,821 trên 39 câu hỏi.
Kết quả cho thấy biểu diễn hàng và ô là cần thiết để tạo bằng chứng chi tiết,
nhưng Cell Match@5 còn 0,487 nên thứ hạng của cell chunk vẫn cần được cải thiện.
Tập 39 câu có độ khó cao hơn vì bao phủ nhiều bảng, nhiều dạng truy vấn và các
câu tra cứu theo khoảng điều kiện. Vì vậy, kết luận RQ5 được giới hạn trong phạm
vi tài liệu QCDT và các bảng đã gán nhãn.

\subsection{Truy xuất thông tin}
\label{subsection:retrieval-results}

BM25, dense, hybrid và hybrid có sắp xếp lại được so sánh trên ba miền dữ liệu.
Các tập SciFact và QASPER được dùng để kiểm tra khả năng chuyển miền, không
phải cơ sở cho tuyên bố đạt kết quả tốt nhất hiện có
\cite{thakur2021beir,wadden2020scifact,dasigi2021qasper}.

\begin{table}[H]
\centering
\scriptsize
\renewcommand{\arraystretch}{1.12}
\caption{Kết quả truy xuất trên SciFact, QCDT và QASPER}
\label{tab:retrieval-master-results}
\begin{tabularx}{\textwidth}{
    |>{\raggedright\arraybackslash}p{0.14\textwidth}
    |>{\raggedright\arraybackslash}p{0.19\textwidth}
    |>{\centering\arraybackslash}p{0.09\textwidth}
    |>{\centering\arraybackslash}p{0.12\textwidth}
    |>{\centering\arraybackslash}p{0.12\textwidth}
    |>{\centering\arraybackslash}X|
}
\hline
\textbf{Tập} & \textbf{Chiến lược} & \textbf{Hit@k} &
\textbf{Recall@k} & \textbf{MRR@k} & \textbf{NDCG@k} \\
\hline
SciFact, \(k=5\) & BM25 & 0,713 & 0,693 & 0,569 & 0,594 \\
& Dense & 0,713 & 0,697 & 0,580 & 0,604 \\
& \textbf{Hybrid} & \textbf{0,793} & \textbf{0,771} & 0,654 & \textbf{0,675} \\
& Hybrid + rerank & 0,780 & 0,762 & \textbf{0,654} & 0,674 \\
\hline
QCDT, \(k=10\) & BM25 & 0,550 & 0,550 & 0,387 & 0,469 \\
& Dense & 0,400 & 0,400 & 0,258 & 0,331 \\
& \textbf{Hybrid} & \textbf{0,600} & \textbf{0,600} & 0,407 & \textbf{0,492} \\
& Hybrid + rerank & 0,575 & 0,575 & \textbf{0,413} & 0,490 \\
\hline
QASPER, \(k=20\) & BM25 & 0,570 & 0,502 & \textbf{0,290} & \textbf{0,717} \\
& Dense & 0,420 & 0,370 & 0,170 & 0,571 \\
& Hybrid & 0,570 & \textbf{0,518} & 0,250 & 0,670 \\
& Hybrid + rerank & 0,570 & \textbf{0,518} & 0,281 & 0,691 \\
\hline
\end{tabularx}
\end{table}

Hybrid đạt Hit@k cao nhất trên SciFact và QCDT, đồng thời tăng Recall@20 trên
QASPER. Tuy nhiên, reranking không cải thiện ổn định và BM25 vẫn có MRR,
NDCG tốt hơn trên QASPER. Trên QCDT, dense-only thấp hơn BM25, phù hợp với
đặc trưng thuật ngữ pháp quy cần khớp từ khóa chính xác. Kết quả RQ3 vì vậy là:
hybrid hữu ích nhưng không chiếm ưu thế trên mọi độ đo và mọi miền.

\section{Thực nghiệm đầu cuối}
\label{section:end-to-end-results}

\subsection{Tập đánh giá chính}

Cấu hình C0 dùng một lần truy xuất theo kế hoạch đã chọn, sau đó kiểm tra bằng
chứng, sinh câu trả lời và tạo dẫn chứng. Cấu hình này không tự động truy xuất
lại khi bộ kiểm tra trả về \texttt{expand\_retrieval} hoặc
\texttt{switch\_strategy}.

\begin{table}[H]
\centering
\small
\renewcommand{\arraystretch}{1.2}
\caption{Kết quả hỏi đáp đầu cuối trên tập đánh giá chính}
\label{tab:primary-e2e-results}
\begin{tabularx}{\textwidth}{
    |>{\raggedright\arraybackslash}p{0.22\textwidth}
    |>{\centering\arraybackslash}p{0.13\textwidth}
    |>{\centering\arraybackslash}p{0.16\textwidth}
    |>{\centering\arraybackslash}p{0.16\textwidth}
    |>{\centering\arraybackslash}X|
}
\hline
\textbf{Tập} & \textbf{Answer Match} & \textbf{Evidence Match} &
\textbf{Grounded Rate} & \textbf{Hallucination Rate} \\
\hline
QCDT, 40 câu & 0,725 & 1,000 & 1,000 & 0,000 \\
\hline
\end{tabularx}
\end{table}

Trên QCDT, bằng chứng được tìm đúng nhưng Answer Match chỉ đạt 0,725. Phân
tích từng câu cho thấy lỗi chủ yếu nằm ở việc tổng hợp thiếu một phần điều kiện
hoặc chọn câu trả lời ngắn chưa khớp đầy đủ đáp án chuẩn. Hallucination Rate bằng 0,000
chỉ có nghĩa không phát hiện ca thiếu căn cứ theo bộ đánh giá và cấu hình này;
đây không phải bảo đảm loại bỏ hallucination trên mọi tài liệu.

\subsection{Tập đánh giá phụ và giới hạn chuyển miền}

\begin{table}[H]
\centering
\small
\renewcommand{\arraystretch}{1.2}
\caption{Kết quả hỏi đáp trên dữ liệu khoa học}
\label{tab:secondary-e2e-results}
\begin{tabularx}{\textwidth}{
    |>{\raggedright\arraybackslash}p{0.15\textwidth}
    |>{\centering\arraybackslash}p{0.09\textwidth}
    |>{\centering\arraybackslash}p{0.11\textwidth}
    |>{\centering\arraybackslash}p{0.11\textwidth}
    |>{\centering\arraybackslash}p{0.12\textwidth}
    |>{\centering\arraybackslash}p{0.14\textwidth}
    |>{\centering\arraybackslash}X|
}
\hline
\textbf{Tập} & \textbf{Số câu} & \textbf{Answer} & \textbf{Evidence} &
\textbf{Grounded} & \textbf{Hallucination} & \textbf{E2E success} \\
\hline
SciFact & 300 & 0,220 & 0,727 & 1,000 & 0,000 & 0,203 \\
\hline
QASPER & 500 & 0,084 & 0,240 & 1,000 & 0,054 & 0,050 \\
\hline
\end{tabularx}
\end{table}

Grounded Rate cao nhưng Answer Match và E2E success thấp cho thấy câu trả lời
có thể bám vào bằng chứng đã chọn mà vẫn thiếu ý hoặc không khớp đáp án chuẩn.
QASPER thường yêu cầu tổng hợp nhiều đoạn và trả lời tự do, trong khi bộ sinh
hiện tại thiên về câu trả lời ngắn, có kiểm soát. Do đó, SciFact và QASPER được
dùng để chỉ ra giới hạn, không được gộp với kết luận chính trên QCDT.

\section{Ablation và phân tích thành phần}
\label{section:ablation-results}

Bảng~\ref{tab:ablation-summary} tóm tắt các ablation chính. Các kết quả chi tiết
về backend bảng, table-aware chunking và chiến lược truy xuất đã được trình bày
ở các bảng trước để tránh lặp số liệu.

\begin{table}[H]
\centering
\small
\renewcommand{\arraystretch}{1.2}
\caption{Tổng hợp vai trò của các thành phần}
\label{tab:ablation-summary}
\begin{tabularx}{\textwidth}{
    |>{\raggedright\arraybackslash}p{0.24\textwidth}
    |>{\raggedright\arraybackslash}p{0.27\textwidth}
    |>{\raggedright\arraybackslash}X|
}
\hline
\textbf{Thành phần} & \textbf{So sánh} & \textbf{Kết luận trong phạm vi thí nghiệm} \\
\hline
Region routing &
Bật so với tắt trên 31 PDF &
Tăng số block bảng và route trace trên QCDT, nhưng tăng thời gian xử lý; trong
thực nghiệm đầu cuối BM25-only, Hit@10 và Answer Match không tăng. \\
\hline
Hybrid TATR &
So với Default và TATR &
Cải thiện Cell F1, Structure F1 và Text Assignment F1; trên 100 mẫu, Exact CSV
đạt 0,530 và còn yếu ở ô gộp, tiêu đề nhiều hàng. \\
\hline
Hybrid retrieval &
So với BM25 và dense &
Tăng Hit@k trên SciFact và QCDT; không tốt nhất trên mọi độ đo QASPER. \\
\hline
Table-aware chunking &
So với chunk thông thường &
Cải thiện Row Match từ 0,000 lên 0,974 và Cell Match@5 từ 0,000 lên 0,487 trên
39 câu hỏi. \\
\hline
Router, evidence checker, citation &
C0 so với các cấu hình tắt từng thành phần &
Router cải thiện QCDT; citation grounding quyết định Grounded Rate. Evidence
checker chưa làm tăng Answer Match trên QCDT trong lần chạy này. \\
\hline
\end{tabularx}
\end{table}

\begin{table}[H]
\centering
\scriptsize
\renewcommand{\arraystretch}{1.15}
\caption{Ablation của luồng hỏi đáp trên QCDT}
\label{tab:qa-ablation-results}
\begin{tabularx}{\textwidth}{
    |>{\raggedright\arraybackslash}p{0.13\textwidth}
    |>{\raggedright\arraybackslash}p{0.24\textwidth}
    |>{\centering\arraybackslash}p{0.11\textwidth}
    |>{\centering\arraybackslash}p{0.12\textwidth}
    |>{\centering\arraybackslash}p{0.12\textwidth}
    |>{\centering\arraybackslash}X|
}
\hline
\textbf{Tập} & \textbf{Cấu hình} & \textbf{Answer} & \textbf{Evidence} &
\textbf{Grounded} & \textbf{Hallucination} \\
\hline
QCDT & C0: routed grounded & 0,725 & 1,000 & 1,000 & 0,000 \\
& Không có router & 0,675 & 0,950 & 1,000 & 0,000 \\
& Không có evidence checker & 0,750 & 1,000 & 1,000 & 0,000 \\
& Không grounding citation & 0,725 & 1,000 & 0,000 & 1,000 \\
\hline
\end{tabularx}
\end{table}

Trên QCDT, tắt evidence checker làm Answer Match tăng nhẹ từ 0,725 lên 0,750,
nên không thể kết luận evidence checker luôn làm tăng độ chính xác câu trả lời.
Tắt citation grounding không làm thay đổi Answer Match nhưng làm giảm mạnh
Grounded Rate và làm Hallucination Rate tăng theo quy tắc đánh giá. Điều đó
cho thấy độ đúng nội dung và khả năng kiểm chứng là hai tiêu chí khác nhau.

\section{Phân tích lỗi và giới hạn}
\label{section:error-analysis}

\begin{table}[H]
\centering
\small
\renewcommand{\arraystretch}{1.2}
\caption{Các nhóm lỗi chính quan sát được}
\label{tab:error-cases}
\begin{tabularx}{\textwidth}{
    |>{\raggedright\arraybackslash}p{0.20\textwidth}
    |>{\raggedright\arraybackslash}p{0.24\textwidth}
    |>{\raggedright\arraybackslash}X
    |>{\raggedright\arraybackslash}p{0.23\textwidth}|
}
\hline
\textbf{Nhóm lỗi} & \textbf{Biểu hiện} & \textbf{Nguyên nhân quan sát được} &
\textbf{Hướng cải thiện} \\
\hline
Câu hỏi quy chế nhiều điều kiện &
Evidence đúng nhưng câu trả lời thiếu một ý hoặc điều kiện. &
Chunk ngắn và bộ sinh ưu tiên câu trả lời trích xuất ngắn; Answer Match QCDT
chỉ đạt 0,725. &
Truy xuất phân cấp theo mục, tổng hợp nhiều bằng chứng. \\
\hline
Truy xuất dense trên QCDT &
Dense Hit@10 đạt 0,400, thấp hơn BM25 0,550. &
Thuật ngữ pháp quy và tên điều khoản cần khớp từ khóa chính xác. &
Embedding tiếng Việt theo miền và hard-negative phù hợp. \\
\hline
Bảng có ô gộp hoặc tiêu đề nhiều hàng &
Cell F1 cao nhưng Exact CSV của Hybrid TATR chỉ đạt 0,530 trên 100 mẫu. &
Sai khác nhỏ trong tách hàng, gán span và thứ tự ô làm toàn bộ CSV không khớp. &
Tái dựng bằng đồ thị và ràng buộc cấu trúc. \\
\hline
Ô bảng nằm ngoài top-5 &
Cell Match@5 đạt 0,487 nhưng Cell Match@10 đạt 0,821 trên 39 câu hỏi. &
Nhiều cell chunk đã tồn tại nhưng điểm xếp hạng chưa đủ cao. &
Reranker chuyên biệt cho row/column header và giá trị ô. \\
\hline
QASPER nhiều bằng chứng &
Answer Match 0,084 và E2E success 0,050. &
Câu hỏi dài, đáp án tự do và cần kết hợp nhiều đoạn trong bài báo. &
Lập kế hoạch bằng chứng và sinh tóm tắt có kiểm chứng. \\
\hline
OCR tài liệu thực tế &
Token F1 FUNSD 0,826, thấp hơn tập scan tổng hợp 1,000. &
Nhiễu ảnh, phông chữ, bố cục và căn chỉnh bbox phức tạp hơn. &
Thêm tập scan tiếng Việt gán nhãn và truyền độ tin cậy OCR. \\
\hline
\end{tabularx}
\end{table}

\subsection{Các ca lỗi tiêu biểu}

\textbf{Ca QCDT q11.} Câu hỏi yêu cầu phân biệt học phần tiên quyết, học trước
và song hành. Hệ thống tìm được bằng chứng liên quan nhưng câu trả lời chủ yếu
nêu điều kiện của đồ án tốt nghiệp và chỉ giải thích học phần song hành. Đây là
trường hợp Evidence Match đúng nhưng Answer Match sai vì chưa tổng hợp đủ ba
khái niệm.

\textbf{Ca QCDT q17.} Câu hỏi hỏi hậu quả khi người học không nộp học phí đúng
hạn. Hệ thống chỉ trả lại nghĩa vụ ``nộp học phí đầy đủ và đúng thời hạn'', không
trả lời hậu quả. Ca này cho thấy khớp từ khóa có thể đưa về đúng đoạn chủ đề nhưng
không bảo đảm đoạn được chọn chứa quan hệ mà câu hỏi yêu cầu.

\textbf{Ca Table QA real\_table\_q006.} Với câu hỏi về cột ``1--2'', cấu hình
table-aware có cell chunk tìm đúng bảng, hàng và cột nhưng cell chưa nằm trong
top-5. Khi tăng lên top-10, cell đúng được thu hồi. Lỗi nằm ở thứ hạng bằng chứng,
không phải thiếu dữ liệu ô trong corpus.

\textbf{Ca Table QA real\_table\_q037.} Câu hỏi về hàng ``thang 10'' trong bảng
quy đổi điểm học phần vẫn chưa khớp cell ở top-10. Bảng này có header rộng và
một ô chứa nhiều giá trị, nên bộ truy xuất dễ tìm đúng bảng nhưng khó đưa đúng
cell lên thứ hạng cao.

\textbf{Ca QASPER về bảy ngôn ngữ Ấn Độ.} Hệ thống lấy một đoạn cùng bài báo
nhưng đoạn đó mô tả chiến lược pooling, không liệt kê bảy ngôn ngữ được hỏi.
Trường hợp này minh họa giới hạn của truy xuất và tổng hợp bằng chứng trên bài
báo dài.

Các kết quả còn hai giới hạn chính. Thứ nhất, tập cấu trúc bảng mở rộng gồm 100
mẫu PubTables proxy với hộp từ chuẩn hóa, nhưng phân nhóm ``không có đường
kẻ'' được xác định bằng heuristic ảnh; tập Table QA mở rộng mới có 39 câu hỏi
trên một tài liệu QCDT. Thứ hai, Grounded Rate và Hallucination Rate phụ thuộc
quy tắc đánh giá; chúng không
thay thế đánh giá thủ công về tính đúng, đầy đủ và phù hợp ngữ cảnh.

\section{Kết luận chương}
\label{section:experiment-conclusion}

\begin{table}[H]
\centering
\small
\renewcommand{\arraystretch}{1.2}
\caption{Kết luận theo câu hỏi nghiên cứu}
\label{tab:rq-conclusions}
\begin{tabularx}{\textwidth}{
    |>{\raggedright\arraybackslash}p{0.10\textwidth}
    |>{\raggedright\arraybackslash}X
    |>{\raggedright\arraybackslash}p{0.31\textwidth}|
}
\hline
\textbf{RQ} & \textbf{Kết luận} & \textbf{Giới hạn} \\
\hline
RQ1 &
Luồng ingest chính đạt success rate 1,000 trên 31 PDF thử nghiệm; bật region
routing tạo route trace, block bảng và metadata nguồn rõ hơn. &
Region routing tăng khả năng truy vết nhưng chưa chứng minh cải thiện QA ở cấu
hình BM25-only; chưa có tập nhãn thủ công để tính Routing Accuracy theo từng
loại vùng. \\
\hline
RQ2 &
Hybrid TATR cải thiện rõ tái dựng bảng so với hai backend còn lại. &
Exact CSV 0,530 trên 100 mẫu; lỗi còn tập trung ở ô gộp và tiêu đề nhiều hàng. \\
\hline
RQ3 &
Hybrid retrieval tăng Hit@k trên SciFact và QCDT. &
Không tốt nhất trên mọi độ đo QASPER; reranking chưa ổn định. \\
\hline
RQ4 &
Cấu hình chính đạt Answer Match 0,725 trên QCDT, đồng thời giữ Evidence Match
1,000 và Grounded Rate 1,000 theo bộ đánh giá. &
Không bảo đảm loại bỏ hallucination ngoài các tập và quy tắc đã đánh giá. \\
\hline
RQ5 &
Table-aware chunking giúp truy xuất đúng hàng và tạo khả năng truy xuất ô. &
Cell Match@5 còn 0,487; tập mở rộng 39 câu mới bao phủ một tài liệu QCDT. \\
\hline
\end{tabularx}
\end{table}

Nhìn chung, kết quả hỗ trợ lựa chọn thiết kế của đồ án: trích xuất PDF theo vùng,
kết hợp hình học bảng với hộp từ, giữ metadata nguồn và đánh giá câu trả lời cùng
bằng chứng. Đồng thời, kết quả trên QASPER, OCR thực tế và Exact CSV cho thấy
hệ thống vẫn còn hạn chế ở tổng hợp nhiều bằng chứng, xử lý scan và tái dựng
bảng phức tạp. Các kết luận trong chương chỉ áp dụng cho dữ liệu, cấu hình và
các tệp kết quả thực nghiệm đã nêu.

\end{document}
